\label{bkm:Ref36969669}{

A BPatch_frame object represents a stack frame. The getCallStack member function of BPatch_thread returns a vector of

BPatch_frame objects representing the frames currently on the stack.}





\bigskip



{

BPatch_frameType getFrameType\index{getFrameType: : : : : : :!}()}



{

Return the type of the stack frame. Possible types are:}



\begin{center}

\tablefirsthead{}

\tablehead{}

\tabletail{}

\tablelasttail{}

\begin{supertabular}{|m{1.5573599in}|m{3.75456in}|}

\hline

{ Frame Type} &

{ Meaning}\\\hline

{ BPatch_frameNormal} &

{ A normal stack frame.}\\\hline

{ BPatch_frameSignal} &

{ A frame that represents a signal invocation.\\\hline

{ BPatch_frameTrampoline} &

{ A frame the represents a call into instrumentation code.}\\\hline

\end{supertabular}

\end{center}



\bigskip



{

void *getFP\index{getFP: : : : : : :!}()}



{

Return the frame pointer for the stack frame.





\bigskip



{

void *getPC\index{getPC: : : : : : :!}()}



{

Returns the program counter associated with the stack frame.}





\bigskip



{

BPatch_function *findFunction\index{findFunction: : : : : : :!}()}



{

Returns the function associated with the stack frame.}



{

BPatch_thread *getThread()}



{

Returns the thread associated with the stack frame.}



{

BPatch_point *getPoint()}



{

BPatch_point *findPoint()}



{

For stack frames corresponding to inserted instrumentation, returns the instrumentation point where that instrumentation

was inserted. For other frames, returns NULL.}



{

bool isSynthesized()}



{

Returns true if this frame was artificially created, false otherwise.}



